% ===================================================================
%  اللوحة رقم ١٥ — السوق. — الأطعمة.
%  Traduction 2026 d'apres texto/io, controlee sur texto/fr, texto/en
%  et texto/es. Consigne : voir texto/ar/10-lawha-01.tex.
% ===================================================================

%%K t15-apar-1 apar
\VUsaut{4.0mm}
\VUtitre{15.0pt}{60}{اللوحة رقم ١٥}
\VUsaut{3.0mm}

%%K t15-tit-1 sub
\VUpk{11.4pt}{40}{السوق. --- الأطعمة.}
\VUsaut{3.0mm}

%%K t15-01-1 p
1. --- معظم التجّار الذين يمدّوننا بما نحتاج من طعام مجتمعون تحت
\VUgras{سقيفة} \textsuperscript{(1)} \VUgras{السوق المسقوفة}، أو في الشوارع
المجاورة.

%%K t15-02-1 p
2. --- \VUgras{بائع لحم الخنزير} \textsuperscript{(2)}، الذي يبيع \VUgras{الفخذ
المقدّد} \textsuperscript{(3)} و\VUgras{النقانق السوداء} \textsuperscript{(4)}
و\VUgras{السجق} \textsuperscript{(5)} و\VUgras{المصارين المحشوّة} \textsuperscript{(6)}
و\VUgras{الشحم} \textsuperscript{(7)}، يقف إلى جانب \VUgras{بائعة الزبدة والجبن}
\textsuperscript{(8)} التي فُرش بسطتها بـ\VUgras{أجبان هولندية} \textsuperscript{(9)}
حمراء القشرة، و\VUgras{كامامبير} \textsuperscript{(10)}، و\VUgras{أقراص غرويير}
\textsuperscript{(11)} كبيرة، و\VUgras{قوالب زبدة} \textsuperscript{(12)} طازجة.

%%K t15-03-1 p
3. --- وأمامها تعرض \VUgras{بائعة خضر} \textsuperscript{(13)} على زبائنها
\VUgras{طماطم} \textsuperscript{(14)} جميلة، و\VUgras{قرنبيطًا} \textsuperscript{(15)}،
و\VUgras{خسًّا} \textsuperscript{(16)}، و\VUgras{ملفوفًا} \textsuperscript{(17)}،
و\VUgras{يقطينًا} \textsuperscript{(19)}، و\VUgras{شمّامًا} \textsuperscript{(18)}، بل
و\VUgras{فطرًا} \textsuperscript{(20)}. لكن \VUgras{عامل جعة} أوقف
\VUgras{عربة توصيله} \textsuperscript{(21)} أمام بسطتها تمامًا، وهي محمّلة
بـ\VUgras{براميل الجعة} \textsuperscript{(22)}، فمنع الزبائن من الاقتراب.
وبستانيّ يجلب لها سلة من \VUgras{الخرشوف} \textsuperscript{(23)}.

%%K t15-tit-2 sub
\VUsaut{4.0mm}
\VUpk{10.2pt}{40}{البيع والشراء.}
\VUsaut{3.0mm}

%%K t15-04-1 p
4. --- والحركة في السوق شديدة، فقد حانت ساعة \VUgras{المزاد} \textsuperscript{(24)}.
و\VUgras{ربّات البيوت} \textsuperscript{(26)} اللواتي يأتين للتسوّق يقفن في
طريقهن أمام بسطة \VUgras{بائعة الأحشاء} \textsuperscript{(25)} ليشترين
\VUgras{كرشًا} أو \VUgras{رأس عجل}.

%%K t15-05-1 p
5. --- و\VUgras{طاهية} \textsuperscript{(27)} بدينة تساوم على \VUgras{تونة} جميلة
جدًا \VUgras{بائعةَ السمك} \textsuperscript{(28)} التي ترفع أسعارها كما يحلو
لها. وقد وصلتها كمية كبيرة من \VUgras{الأنقليس والسمك موسى والتروتة}
و\VUgras{الشبّوط}. و\VUgras{قدّها} و\VUgras{بلحها البحري} طازجان جدًا.

%%K t15-tit-3 sub
\VUsaut{4.0mm}
\VUpk{10.2pt}{40}{الرخيص والغالي.}
\VUsaut{3.0mm}

%%K t15-06-1 p
6. --- و\VUgras{بائعة الدواجن} \textsuperscript{(29)} المستقرة إلى جانبها أمينة
جدًا. فإذا باعت \VUgras{ديكًا روميًا} أو \VUgras{إوزة} أو \VUgras{بطة} أو
\VUgras{دجاجة} عادية، لم تطلب أكثر من السعر العادل.

%%K t15-07-1 p
7. --- وفي زاوية الساحة، في الطابق الأول، استقرت منذ قريب \VUgras{بائعة
بياضات} \textsuperscript{(30)}، والمشترون واثقون دائمًا أن يجدوا عندها اختيارًا
حسنًا جدًا من \VUgras{المفارش} و\VUgras{المناديل والمآزر} و\VUgras{فوط
المطبخ}. وفي الطابق نفسه أقام \VUgras{سكّاكيني} \textsuperscript{(31)} ورشته.
يسنّ \VUgras{سكاكين} بائعات السوق ويصنع للبستانيين \VUgras{مناجل تقليم} من
أصلب \VUgras{فولاذ}.

%%K t15-tit-4 sub
\VUsaut{4.0mm}
\VUpk{10.2pt}{40}{دكان البقالة.}
\VUsaut{3.0mm}

%%K t15-08-1 p
8. --- وتحت ورشته فُتح منذ غير بعيد دكان أطعمة كبير. لم يعد
\VUgras{دكان البقالة} \textsuperscript{(32)} البسيط المتواضع في الزمن الغابر،
الذي لم يكن يبيع إلا الأصناف المعتادة --- \VUgras{الشاي} \textsuperscript{(33)}
و\VUgras{الشوكولاتة} \textsuperscript{(34)} و\VUgras{قالب السكر} \textsuperscript{(37)}
و\VUgras{الصابون} \textsuperscript{(38)}. هنا يُباع كل شيء:
\VUgras{بسكويت مقرمش}، و\VUgras{معلّبات} \textsuperscript{(35)}،
و\VUgras{مشروبات} \textsuperscript{(36)}، و\VUgras{رم وبراندي وكيرش ويانسون}
و\VUgras{كوراساو}. ويوجد الصيد هنا بالسهولة نفسها --- \VUgras{أرنب بري}
\textsuperscript{(39)}، و\VUgras{سُمنة} \textsuperscript{(40)}، و\VUgras{حجل}
\textsuperscript{(41)}، و\VUgras{سُمّان} \textsuperscript{(42)}، و\VUgras{بط بري}
\textsuperscript{(43)}، أو \VUgras{تدرج} \textsuperscript{(44)} --- كما توجد فاكهة
المناطق الحارة مثل \VUgras{الموز} \textsuperscript{(46)} و\VUgras{الأناناس}.

%%K t15-09-1 p
9. --- و\VUgras{مساعد} \textsuperscript{(45)} بقّال شديد اللطف مستعد أن يناول ما
يُطلب منه: \VUgras{مربّى} \textsuperscript{(47)} من الكشمش أو السفرجل،
و\VUgras{شحمًا} \textsuperscript{(48)}، و\VUgras{مخللًا} \textsuperscript{(49)}، أو
\VUgras{سردينًا} \textsuperscript{(50)} مملّحًا.

%%K t15-09-2 p
وذلك مريح جدًا لـ\VUgras{الزبائن} \textsuperscript{(51)}، فهم يجدون، إلى جانب
أشد الأصناف اعتيادًا، أندر البواكير وألذّ الفاكهة: \VUgras{كمثرى}
\textsuperscript{(52)} عصيرية، و\VUgras{برقوقًا} \textsuperscript{(53)}،
و\VUgras{عنبًا} \textsuperscript{(54)}، و\VUgras{خوخًا} \textsuperscript{(55)}،
و\VUgras{لوزًا} \textsuperscript{(56)}، و\VUgras{جوزًا} \textsuperscript{(57)}.

%%K t15-tit-5 sub
\VUsaut{4.0mm}
\VUpk{10.2pt}{40}{الزبائن.}
\VUsaut{3.0mm}

%%K t15-10-1 p
10. --- و\VUgras{الأرز} \textsuperscript{(58)} و\VUgras{العدس} \textsuperscript{(59)}
إلى جانب علبة \VUgras{الرنكة} \textsuperscript{(60)} المدخّنة؛
و\VUgras{البطاطا} \textsuperscript{(61)} و\VUgras{الفاصولياء} \textsuperscript{(63)}
تجاور \VUgras{التفاح} \textsuperscript{(62)} و\VUgras{التين} \textsuperscript{(64)}
و\VUgras{القراصيا} \textsuperscript{(65)} و\VUgras{البندق} \textsuperscript{(66)}.
ويوجد في هذا الدكان أيضًا \VUgras{بيض} \textsuperscript{(67)} طازج
و\VUgras{زيتون} \textsuperscript{(68)}، فتمتلئ \VUgras{السلال} \textsuperscript{(70)}
و\VUgras{شِباك التسوق} \textsuperscript{(73)} بسرعة كبيرة.

%%K t15-11-1 p
11. --- و\VUgras{قهوة} \textsuperscript{(72)} هذا البيت الممتاز نالت شهرة تستحقها
بين \VUgras{الخادمات} \textsuperscript{(69)} والطاهيات، فـ\VUgras{البقّال}
\textsuperscript{(71)} العجوز يحمّصها بنفسه بأشد العناية.

%%K t15-tit-6 sub
\VUsaut{4.0mm}
\VUpk{10.2pt}{40}{أشياء تُرى.}
\VUsaut{3.0mm}

%%K t15-12-1 p
12. --- وليس كل من يأتي السوق يأتيه للشراء. ففي الرواح والمجيء الطريف في
الزقاق المؤدي إلى السقيفة، يُرى أشد الناس تنوعًا. فإلى جانب \VUgras{عامل
مسلخ} \textsuperscript{(75)} طيّب النفس يدفع أمامه \VUgras{عربة يد} \textsuperscript{(74)}،
ها هما جنديان في إجازة، فخوران جدًا بعرض زيّيهما الزاهيين:
\VUgras{الفارس} \textsuperscript{(76)} في \VUgras{سترته} الزرقاء و\VUgras{قلنسوته
ذات الريش}، و\VUgras{عريف الزواف} في سرواله الفضفاض وعلى رأسه
\VUgras{طربوش}.

%%K t15-13-1 p
13. --- و\VUgras{الخبّاز} \textsuperscript{(80)}، واقفًا على عتبة \VUgras{المخبز}
\textsuperscript{(78)}، قد عجن لتوّه عجين \VUgras{خبزه} \textsuperscript{(79)} وأخرج
من الفرن \VUgras{الكرواسان} \textsuperscript{(81)} و\VUgras{الأرغفة الصغيرة}
\textsuperscript{(82)} و\VUgras{الأرغفة المستديرة} \textsuperscript{(83)} التي في
الواجهة.

%%K t15-14-1 p
14. --- وفوق المخبز تسكن \VUgras{خيّاطة} \textsuperscript{(84)} ماهرة تستخدم عدة
\VUgras{مطرّزات} \textsuperscript{(85)}. وإلى جانبها تسكن \VUgras{غسّالة وكاوية}
\textsuperscript{(86)} تنشّي \VUgras{البياضات} \textsuperscript{(88)} بعد غسلها
وتجفيفها، وتكويها بـ\VUgras{مِكواة} \textsuperscript{(87)}.

%%K t15-tit-7 sub
\VUsaut{4.0mm}
\VUpk{10.2pt}{40}{اللحم والسمك والخضر.}
\VUsaut{3.0mm}

%%K t15-15-1 p
15. --- وفي \VUgras{الملحمة} \textsuperscript{(89)} في الطابق الأرضي يُباع اللحم
على اختلافه --- \VUgras{ضأن} \textsuperscript{(90)} أو \VUgras{بقر} \textsuperscript{(94)}
أو \VUgras{عجل} \textsuperscript{(92)} --- في صورة \VUgras{أفخاذ وشرائح وقطع شواء}
و\VUgras{ريش}. و\VUgras{صبي اللحّام} \textsuperscript{(93)} يقطّع هذا اللحم كله
ويعزل \VUgras{عظام النخاع} \textsuperscript{(96)}. وعلى باب الملحمة تقف
\VUgras{بائعة محار} \textsuperscript{(97)} تبيع \VUgras{المحار} \textsuperscript{(98)}.
وهي تفتحه إن شئت.

%%K t15-16-1 p
16. --- وهؤلاء التجّار كلهم يدفعون رخصة أو رسم بسطة؛ ولذلك يطارد
\VUgras{الشرطي} \textsuperscript{(100)} بلا رحمة \VUgras{بائعات الخضر المتجولات}
\textsuperscript{(99)} المسكينات اللواتي ينافسنهم بحملهن على عرباتهن
\VUgras{جزرًا وفجلًا ولفتًا وكرّاثًا} أو \VUgras{حزم هليون وبازلاء طازجة
وسبانخ وحمّاض ورشاد} و\VUgras{بقدونس}.

%%K t15-tit-8 sub
\VUsaut{4.0mm}
\VUpk{10.2pt}{40}{احذر الشرطي!}
\VUsaut{3.0mm}

%%K t15-17-1 p
17. --- والصبي الذي يبيع \VUgras{الثوم} \textsuperscript{(101)} و\VUgras{البصل}
يعلم جيدًا أنه هو أيضًا لا يُسمح له ببيع بضاعته قرب السوق. فيفرّ مخاطرًا
بقلب سلة \VUgras{السرطانات} و\VUgras{جراد البحر الصغير}
و\VUgras{القريدس} التي لـ\VUgras{بائعة} \textsuperscript{(102)}
\VUgras{الجراد البحري الشوكي} و\VUgras{الكركند}.

%%K t15-18-1 p
18. --- والجميع يسعون ويحدثون ضجيجًا كبيرًا؛ لكن ذلك لا يمنع من أن يُسمع
بوضوح نداء \VUgras{الزجّاج} \textsuperscript{(103)} الجوّال، ولا صيحة
\VUgras{الفحّام} \textsuperscript{(104)} الذي ترك عربته لحظة ليوصل \VUgras{كيس
فحم} \textsuperscript{(105)}.
\VUsaut{6.0mm}
\VUfilet{22mm}
